\chapter{État de l'art}
\label{chap:etat}

Concevoir un langage et l'outil capable de l'exécuter n'est pas une démarche
isolée. C'est au contraire s'inscrire dans une tradition technique et théorique
qui s'est constituée sur plus d'un demi-siècle, et qui a fini par dégager des
méthodes et des outils devenus standards. Avant de présenter mes propres choix,
il me paraît donc indispensable de situer le projet dans ce paysage : de rappeler
d'où viennent les concepts que j'ai utilisés, quelles solutions existent pour
chaque problème que j'ai eu à résoudre, et pourquoi j'ai retenu telle approche
plutôt que telle autre. C'est l'objet de ce chapitre, organisé en une revue des
travaux et des outils existants, suivie du positionnement de mon travail.

\section{Revue de littérature}

\subsection{Travaux fondateurs}

Tout ce projet repose, souvent sans qu'on y pense, sur des résultats théoriques
établis entre les années 1950 et 1970. Le premier de ces apports est la
classification des grammaires formelles proposée par Noam Chomsky en
1956~\cite{chomsky}. Chomsky a montré qu'on peut hiérarchiser les langages selon
la puissance de la machine nécessaire pour les reconnaître. Deux niveaux de cette
hiérarchie nous intéressent directement. Les langages dits réguliers, les plus
simples, sont reconnus par des automates finis : ils suffisent à décrire la forme
des mots du langage, c'est-à-dire les nombres, les identifiants ou les
opérateurs. Les langages dits hors-contexte, plus expressifs, sont reconnus par
des automates à pile : ils sont nécessaires pour décrire l'imbrication des
parenthèses, des blocs et des expressions. Cette distinction n'est pas une
subtilité académique. Elle se traduit très concrètement dans l'architecture de
n'importe quel compilateur, et dans la mienne en particulier, par la séparation
en deux étages bien distincts. L'analyse lexicale, confiée à Flex, traite le
niveau régulier. L'analyse syntaxique, confiée à Bison, traite le niveau
hors-contexte. J'ai trouvé éclairant de constater que cette frontière théorique
correspond exactement à la frontière entre mes deux premiers fichiers source.

Le deuxième apport fondateur concerne l'analyse syntaxique elle-même. En 1965,
Donald Knuth a décrit l'analyse LR~\cite{knuth}, une méthode ascendante,
déterministe et efficace, capable de reconnaître une très large classe de
grammaires hors-contexte. Le mot « ascendante » signifie que l'analyseur part des
symboles élémentaires et reconstruit progressivement les structures
englobantes, par opposition aux méthodes descendantes qui partent du but et le
décomposent. Les variantes plus économes en mémoire de cette méthode, en
particulier l'analyse LALR(1) qui n'utilise qu'un seul symbole de prévision, sont
précisément celles que met en \oe uvre l'outil Bison. Autrement dit, lorsque
j'écris une grammaire pour Bison, je m'appuie sans le réécrire sur un algorithme
qui a près de soixante ans et qui fait toujours référence.

L'ensemble de cette chaîne, de l'analyse lexicale jusqu'à la génération de code en
passant par les arbres syntaxiques et les optimisations, a été synthétisé et
codifié dans un ouvrage devenu la référence absolue du domaine, communément appelé
le « Dragon Book », d'Aho, Lam, Sethi et Ullman~\cite{dragon}. La structure même
de mon projet, et donc de ce rapport, suit fidèlement le découpage proposé par cet
ouvrage. Je m'y suis appuyé en particulier pour comprendre le rôle pivot de
l'arbre syntaxique abstrait et la notion de table des symboles.

\subsection{Des outils générateurs plutôt que du code écrit à la main}

Sur le plan des outils, l'apport décisif date de 1975, aux laboratoires Bell, avec
deux programmes qui ont changé la manière d'écrire les compilateurs. Le premier,
\texttt{lex}, dû à Mike Lesk et Eric Schmidt~\cite{lesk}, engendre un analyseur
lexical à partir d'une simple liste d'expressions régulières. Le second,
\texttt{yacc} pour \emph{Yet Another Compiler-Compiler}, dû à Stephen
Johnson~\cite{johnson}, engendre un analyseur syntaxique à partir d'une grammaire.
L'idée commune à ces deux outils est qu'on ne devrait pas écrire un analyseur à la
main, instruction par instruction, mais le \emph{décrire} de façon déclarative et
laisser une machine produire le code correspondant. Cette idée s'est imposée
durablement, et les réimplémentations libres de ces outils, Flex~\cite{flex} pour
l'analyse lexicale et Bison~\cite{bison} pour l'analyse syntaxique, sont
aujourd'hui omniprésentes. Ce sont elles que j'ai employées. L'ouvrage de John
Levine~\cite{levine} m'a servi de guide pratique pour leur utilisation conjointe,
notamment pour comprendre comment le lexer et le parseur communiquent à travers la
variable \texttt{yylval} et la table des tokens.

\subsection{Les stratégies d'exécution}

Une fois le programme analysé et représenté en mémoire, encore faut-il l'exécuter.
Plusieurs stratégies existent, et le choix de l'une d'elles est probablement la
décision d'architecture la plus structurante d'un tel projet. Je les présente ici
de la plus simple à la plus performante, car cette progression correspond aussi à
une montée en complexité de réalisation.

La première stratégie est l'interprétation par parcours d'arbre, souvent désignée
par son nom anglais de \emph{tree-walking}. Le principe consiste à parcourir
récursivement l'arbre syntaxique et à évaluer chaque n\oe ud au moment où on le
rencontre. Un n\oe ud d'addition évalue ses deux fils puis additionne les
résultats, un n\oe ud de condition évalue son test puis descend dans la bonne
branche, et ainsi de suite. C'est l'approche la plus directe et, surtout, la plus
lisible : la sémantique du langage se lit presque mot pour mot dans le code de
l'évaluateur. Son inconvénient est la lenteur, car chaque exécution d'une
instruction suppose de re-parcourir une portion de l'arbre et de déterminer à
chaque n\oe ud de quel type il s'agit.

La deuxième stratégie repose sur une machine virtuelle à bytecode. Le programme
est d'abord traduit dans un jeu d'instructions compact et de bas niveau, le
bytecode, qui est ensuite exécuté par une boucle d'interprétation dédiée. C'est
l'approche retenue par des langages très répandus comme Python, Java ou Lua. Elle
est sensiblement plus rapide que le parcours d'arbre, car le travail d'analyse des
structures est fait une fois pour toutes au moment de la traduction, tout en
restant portable d'une machine à l'autre.

La troisième stratégie, la compilation à la volée ou \emph{JIT}, va plus loin en
traduisant en code machine, pendant l'exécution, les portions de programme les
plus sollicitées. C'est l'approche des moteurs JavaScript modernes ou de la
machine virtuelle Java dans sa version optimisée. Elle offre d'excellentes
performances au prix d'une complexité d'implémentation considérable. La quatrième
et dernière stratégie est la compilation anticipée vers du code machine, telle que
la pratiquent GCC ou Clang pour le langage C. Elle produit les binaires les plus
rapides, mais elle suppose une étape de compilation séparée et un résultat lié à
une architecture précise.

\subsection{Analyse comparative}

Le tableau~\ref{tab:approches} résume les compromis que présentent ces quatre
stratégies. Il ne faut pas le lire comme un classement, car aucune approche n'est
supérieure dans l'absolu : tout dépend de l'objectif poursuivi. Pour un produit
industriel soucieux de performance, le bytecode ou la compilation s'imposent. Pour
un projet dont la finalité est de \emph{comprendre} et de \emph{montrer} comment
un langage fonctionne, c'est au contraire la lisibilité qui prime, et le parcours
d'arbre devient le meilleur choix.

\begin{table}[ht]
\centering
\caption{Comparaison des stratégies d'exécution.}
\label{tab:approches}
\begin{tabular}{@{}lccc@{}}
\toprule
\textbf{Approche} & \textbf{Performance} & \textbf{Complexité} & \textbf{Lisibilité} \\
\midrule
Parcours d'arbre        & faible      & faible   & excellente \\
Machine à bytecode      & moyenne     & moyenne  & moyenne \\
Compilation JIT         & élevée      & élevée   & faible \\
Compilation anticipée   & très élevée & élevée   & faible \\
\bottomrule
\end{tabular}
\end{table}

Un point mérite d'être souligné, car il a orienté toute la suite du projet.
\ilang{} illustre en réalité deux points opposés de cet axe. La version
principale, l'interpréteur, relève du parcours d'arbre. L'extension que j'ai
développée par curiosité, le générateur de code assembleur, relève de la
compilation anticipée. Or les deux partent du \emph{même} arbre syntaxique. Cette
coïncidence n'en est pas une : elle matérialise le fait, central dans toute la
théorie de la compilation, que l'analyse d'un programme est indépendante de ce
qu'on décide d'en faire ensuite. C'est cette idée que j'ai trouvée la plus
instructive de tout le projet, et j'y reviendrai longuement.

\subsection{Les familles d'outils d'analyse syntaxique}

Puisque le c\oe ur de mon travail repose sur un générateur d'analyseur syntaxique,
il est utile de rappeler les grandes familles d'outils disponibles, afin de mieux
justifier ensuite le choix de Bison. La première famille regroupe les générateurs
ascendants de type LR ou LALR, dont Yacc et Bison sont les représentants les plus
connus. Ils construisent automatiquement une table d'analyse à partir de la
grammaire, reconnaissent une large classe de langages, et possèdent une vertu
précieuse pour qui apprend : ils signalent les ambiguïtés de la grammaire sous
forme de conflits. La deuxième famille regroupe les générateurs descendants de
type LL, comme ANTLR ou JavaCC. Ils sont souvent jugés plus intuitifs à déboguer,
car le déroulement de l'analyse suit visiblement la structure de la grammaire,
mais ils sont restreints à une classe de grammaires plus étroite et ont
longtemps été gênés par la récursion à gauche, qu'il faut éliminer à la main.

La troisième famille n'est pas un outil mais une technique : la descente récursive
écrite directement à la main, sans générateur. Elle est souple, lisible pour de
petits langages, et ne dépend d'aucun outil externe. Son défaut majeur, du point
de vue de ce projet, est que la grammaire y reste implicite, dispersée dans le
code, et que rien ne signale ses éventuelles ambiguïtés. La quatrième famille,
plus récente, regroupe les grammaires d'expression d'analyse, dites PEG, et les
combinateurs d'analyseurs. Elles fusionnent souvent l'analyse lexicale et
syntaxique et reposent sur un choix ordonné des règles, ce qui présente l'avantage
de ne jamais produire d'ambiguïté, mais au prix de la masquer plutôt que de la
révéler. Cette dernière propriété est exactement à l'opposé de ce que je
recherchais : je voulais un outil qui me \emph{dise} si ma grammaire était
ambiguë, pas un outil qui décide silencieusement à ma place.

\section{Positionnement}

\subsection{Lacunes de l'approche antérieure}

Le cahier des charges fait explicitement référence à un exercice antérieur du
cours, le mini-langage \textsc{Parlot}, qui réalisait son analyse lexicale au
moyen d'expressions régulières écrites en Python et son analyse syntaxique par une
descente récursive codée à la main. Cette approche fonctionne, et elle a le mérite
de la simplicité initiale, mais elle souffre de deux faiblesses que j'ai cherché à
corriger. La première est que la grammaire n'y est formalisée nulle part. Elle est
diluée dans le code impératif de la descente récursive, si bien qu'on ne peut ni
la lire d'un seul tenant, ni raisonner sur ses propriétés, ni détecter ses
ambiguïtés. La seconde faiblesse découle de la première : comme aucune description
formelle n'existe, rien ne vérifie que tous les cas sont traités. C'est au
programmeur seul d'y veiller, et la moindre ambiguïté passe inaperçue jusqu'à ce
qu'elle provoque un comportement inattendu.

\subsection{Opportunités et justification de mon approche}

Mon projet corrige ces deux lacunes en adoptant les outils standards du cours. En
décrivant la grammaire de façon déclarative dans le fichier Bison, j'obtiens deux
bénéfices que je considère comme décisifs. Le premier est que la grammaire devient
un objet explicite, écrit en notation BNF, que l'on peut lire, discuter, vérifier
et faire évoluer indépendamment du reste du code. Le second, plus important
encore, est que Bison détecte automatiquement les ambiguïtés sous forme de
conflits. Cette détection n'est pas restée théorique dans mon cas : elle m'a
permis de vérifier formellement que ma grammaire était non ambiguë, et même de
faire une découverte qui m'a marqué, l'absence du fameux conflit du « dangling
else », que je détaille au chapitre~\ref{chap:discussion}.

Reste à justifier le choix de l'interprétation par parcours d'arbre plutôt qu'une
machine à bytecode. Ce choix obéit au même principe directeur que le précédent, à
savoir privilégier la clarté. Une machine à bytecode aurait été plus rapide, mais
elle aurait introduit un niveau d'indirection supplémentaire, le jeu
d'instructions intermédiaire, qu'il aurait fallu concevoir, documenter et déboguer
pour un gain de performance sans intérêt dans un cadre pédagogique. Le parcours
d'arbre, au contraire, maintient une correspondance immédiate et visible entre la
structure du programme, représentée par l'arbre, et sa signification, codée dans
l'évaluateur. Ce choix s'est révélé pleinement justifié a posteriori : sa
régularité m'a permis d'ajouter trois extensions au langage, puis un générateur de
code complet, sans jamais avoir à remettre en cause l'architecture de départ.
